Uns bleibt jetzt nur noch, eine allgemeine Methode\SourceLeaf{246}%
 für die oben mit (I) bezeichnete Transformation anzugeben. Welche Funktionen von $x,x',x''$ usw. auch für $y,y',y''$ usw. angenommen werden, so ist klar, dass sich, wenigstens allgemein gesprochen, durch Elimination die Größen $x',x'',x'''$ usw. als Funktionen von $x,y,y',y''$ usw. darstellen lassen, deren Differentiation $n-1$ Gleichungen ergibt:
\[
\begin{aligned}
\mathrm{d}x'&=\xi'\mathrm{d}x+\alpha'\mathrm{d}y+\beta'\mathrm{d}y'+\gamma'\mathrm{d}y''+\text{usw.}\\
\mathrm{d}x''&=\xi''\mathrm{d}x+\alpha''\mathrm{d}y+\beta''\mathrm{d}y'+\gamma''\mathrm{d}y''+\text{usw.}\\
\mathrm{d}x'''&=\xi'''\mathrm{d}x+\alpha'''\mathrm{d}y+\beta'''\mathrm{d}y'+\gamma'''\mathrm{d}y''+\text{usw.}\quad\text{u.\,s.\,w.}
\end{aligned}
\]
Hier sind also die Koeffizienten $\xi',\xi'',\xi'''$ usw., $\alpha',\alpha'',\alpha'''$ usw. Funktionen von $x,y,y',y''$ usw. In dieser Hinsicht werden wir ihre partiellen Ableitungen nach $x$ durch Einschließen in Klammern kennzeichnen. Offenbar können jene Größen auch als Funktionen von $x,x',x'',x'''$ usw. betrachtet werden; in dieser Hinsicht wollen wir die partiellen Ableitungen nach $x$ ohne Klammern schreiben, sodass $\left(\frac{\mathrm{d}\xi'}{\mathrm{d}x}\right)$ sorgfältig von $\frac{\mathrm{d}\xi'}{\mathrm{d}x}$ unterschieden werden muss. Dasselbe gilt für $\left(\frac{\mathrm{d}p}{\mathrm{d}x}\right)$ und $\frac{\mathrm{d}p}{\mathrm{d}x}$ usw. Damit nun $\Omega$ nach Einsetzen jener Werte von $\mathrm{d}x',\mathrm{d}x'',\mathrm{d}x'''$ usw. die vorgeschriebene Form erhält, muss offenbar \emph{erstens} $\mathrm{d}x$ wegfallen, also folgende Bedingungsgleichung [1] gelten:
\[
0=p+p'\xi'+p''\xi''+p'''\xi'''+\text{usw.}
\]
Ferner sollen die Koeffizienten von $\mathrm{d}y,\mathrm{d}y',\mathrm{d}y''$ usw., nämlich
\[
\begin{aligned}
p'\alpha'+p''\alpha''+p'''\alpha'''+\text{usw.}&=A\\
p'\beta'+p''\beta''+p'''\beta'''+\text{usw.}&=B\\
p'\gamma'+p''\gamma''+p'''\gamma'''+\text{usw.}&=C,\quad\text{u.\,s.\,w.}
\end{aligned}
\]
die Form $\lambda q,\lambda q',\lambda q''$ usw. erhalten, sodass $q,q',q''$ usw. nur Funktionen von $y,y',y''$ usw. werden. Damit dies geschieht, müssen wir \emph{zweitens} fordern [2]:
\[
\frac{1}{A}\cdot\left(\frac{\mathrm{d}A}{\mathrm{d}x}\right)
=\frac{1}{B}\cdot\left(\frac{\mathrm{d}B}{\mathrm{d}x}\right)
=\frac{1}{C}\cdot\left(\frac{\mathrm{d}C}{\mathrm{d}x}\right)\text{usw.}
=\frac{1}{\lambda}\cdot\left(\frac{\mathrm{d}\lambda}{\mathrm{d}x}\right)
\]
Nun ist aber
\[
\begin{aligned}
\left(\frac{\mathrm{d}A}{\mathrm{d}x}\right)={}&
p'\left(\frac{\mathrm{d}\alpha'}{\mathrm{d}x}\right)
+p''\left(\frac{\mathrm{d}\alpha''}{\mathrm{d}x}\right)
+p'''\left(\frac{\mathrm{d}\alpha'''}{\mathrm{d}x}\right)+\text{usw.}\\
&+\alpha'\left(\frac{\mathrm{d}p'}{\mathrm{d}x}\right)
+\alpha''\left(\frac{\mathrm{d}p''}{\mathrm{d}x}\right)
+\alpha'''\left(\frac{\mathrm{d}p'''}{\mathrm{d}x}\right)+\text{usw.}
\end{aligned}
\]
\SourceLeaf{247}
Setzt man hier
\[
\left(\frac{\mathrm{d}\alpha'}{\mathrm{d}x}\right)=\frac{\mathrm{d}\xi'}{\mathrm{d}y},\qquad
\left(\frac{\mathrm{d}\alpha''}{\mathrm{d}x}\right)=\frac{\mathrm{d}\xi''}{\mathrm{d}y}.\quad\text{usw.}
\]
ein und subtrahiert
\[
\begin{aligned}
0={}&\frac{\mathrm{d}p}{\mathrm{d}y}
+\xi'\cdot\frac{\mathrm{d}p'}{\mathrm{d}y}
+\xi''\cdot\frac{\mathrm{d}p''}{\mathrm{d}y}+\text{usw.}\\
&+p'\cdot\frac{\mathrm{d}\xi'}{\mathrm{d}y}
+p''\cdot\frac{\mathrm{d}\xi''}{\mathrm{d}y}+\text{usw.}
\end{aligned}
\]
also die Gleichung, die durch Differentiation von [1] nach $y$ entsteht, so ergibt sich
\[
\left(\frac{\mathrm{d}A}{\mathrm{d}x}\right)=
\left\{\begin{aligned}
&\phantom{-\frac{\mathrm{d}p}{\mathrm{d}y}}+\alpha'\cdot\left(\frac{\mathrm{d}p'}{\mathrm{d}x}\right)
+\alpha''\cdot\left(\frac{\mathrm{d}p''}{\mathrm{d}x}\right)+\text{usw.}\\
&-\frac{\mathrm{d}p}{\mathrm{d}y}-\xi'\cdot\frac{\mathrm{d}p'}{\mathrm{d}y}
-\xi''\cdot\frac{\mathrm{d}p''}{\mathrm{d}y}-\text{usw.}
\end{aligned}\right.
\]
\Needspace{22\baselineskip}
Da man ferner
\[
\begin{aligned}
\left(\frac{\mathrm{d}p'}{\mathrm{d}x}\right)&=\frac{\mathrm{d}p'}{\mathrm{d}x}
+\xi'\cdot\frac{\mathrm{d}p'}{\mathrm{d}x'}
+\xi''\cdot\frac{\mathrm{d}p'}{\mathrm{d}x''}+\text{usw.}\\
\left(\frac{\mathrm{d}p''}{\mathrm{d}x}\right)&=\frac{\mathrm{d}p''}{\mathrm{d}x}
+\xi'\cdot\frac{\mathrm{d}p''}{\mathrm{d}x'}
+\xi''\cdot\frac{\mathrm{d}p''}{\mathrm{d}x''}+\text{usw.},\quad\text{u.\,s.\,w.}
\end{aligned}
\]
\[
\begin{aligned}
\frac{\mathrm{d}p}{\mathrm{d}y}&=\alpha'\cdot\frac{\mathrm{d}p}{\mathrm{d}x'}
+\alpha''\cdot\frac{\mathrm{d}p}{\mathrm{d}x''}
+\alpha'''\cdot\frac{\mathrm{d}p}{\mathrm{d}x'''}+\text{usw.}\\
\frac{\mathrm{d}p'}{\mathrm{d}y}&=\alpha'\cdot\frac{\mathrm{d}p'}{\mathrm{d}x'}
+\alpha''\cdot\frac{\mathrm{d}p'}{\mathrm{d}x''}
+\alpha'''\cdot\frac{\mathrm{d}p'}{\mathrm{d}x'''}+\text{usw.},\quad\text{u.\,s.\,w.}
\end{aligned}
\]
hat, erhält man nach diesen Substitutionen
\[
\left(\frac{\mathrm{d}A}{\mathrm{d}x}\right)=k'\alpha'+k''\alpha''+k'''\alpha'''+\text{usw.}
\]
wobei
\[
\begin{array}{r@{}l@{}l@{}l@{}l@{}l}
k'={}&(1,0)&\quad *\quad&+(1,2)\xi''&+(1,3)\xi'''&+\text{usw.}\\
k''={}&(2,0)&+(2,1)\xi'&\quad *\quad&+(2,3)\xi'''&+\text{usw.}\\
k'''={}&(3,0)&+(3,1)\xi'&+(3.2)\xi''-&\quad *\quad&+\text{usw.}
\end{array}
\]
\Needspace{9\baselineskip}
usw. gilt, wenn man der Kürze halber allgemein
\[
\frac{\mathrm{d}p^{(\mu)}}{\mathrm{d}x^{(\nu)}}
-\frac{\mathrm{d}p^{(\nu)}}{\mathrm{d}x^{(\mu)}}
\]
mit $(\mu,\nu)$ bezeichnet, sodass allgemein $(\mu,\mu)=0$ und $(\nu,\mu)=-(\mu.\nu)$ gelten.
Ferner sieht man leicht, dass auch
\[
\begin{aligned}
\left(\frac{\mathrm{d}B}{\mathrm{d}x}\right)&=k'\beta'+k''\beta''+k'''\beta'''+\text{usw.}\\
\left(\frac{\mathrm{d}C}{\mathrm{d}x}\right)&=k'\gamma'+k''\gamma''+k'''\gamma'''+\text{usw.}
\end{aligned}
\]
\SourceLeaf{248}
usw. gilt und daher die Gleichungen [2] erfüllt sind, wenn \(k',\,k'',\,k'''\) usw. den Größen \(p'.\,p'',\,p'''\) usw. jeweils proportional werden. Setzt man außerdem noch
\[
k=\quad *\quad +(0,1)\xi'+(0,2)\xi''+(0,3)\xi'''+\text{usw.}
\]
so hat man die identische Gleichung
\[
0=k+k'\xi'+k''\xi''+k'''\xi'''+\text{usw.}
\]
Aus ihr und [1] folgt leicht, dass auch \(k\) der Größe \(p\) proportional sein muss; diese letzte Proportionalität kann an die Stelle der Gleichung [1] treten. Mithilfe der \(n-1\) Gleichungen
\[
\frac{k}{p}=\frac{k'}{p'}=\frac{k''}{p''}=\frac{k'''}{p'''}=\text{u.\ s.\ w.}
\]
können nun die bisher unbekannten Funktionen \(\xi',\,\xi'',\,\xi'''\) usw., deren Anzahl ebenfalls \(n-1\) ist, bestimmt werden. Eine nähere Betrachtung zeigt jedoch, dass diese Bestimmung nur für gerade Werte von \(n\) ausführbar ist. Bei ungeradem \(n\) fällt stets, sobald alle Größen \(\xi'.\,\xi'',\,\xi'''\) usw. bis auf eine eliminiert sind, auch diese von selbst weg, und es bleibt nur eine \emph{Bedingungsgleichung} zwischen \(p,\,p',\,p''\) usw. übrig. Darin liegt der Grund, weshalb die Transformation (I) auf gerade Werte von \(n\) beschränkt werden muss.

Die Bedingungen der Transformation I sind also jetzt darauf zurückgeführt, dass die partiellen Ableitungen \(\bigl(\frac{\mathrm{d}x'}{\mathrm{d}x}\bigr),\ \bigl(\frac{\mathrm{d}x''}{\mathrm{d}x}\bigr),\ \bigl(\frac{\mathrm{d}x'''}{\mathrm{d}x}\bigr)\) usw., soweit \(x',\,x'',\,x'''\) usw. als Funktionen von \(x,\,y,\,y',\,y''\) usw. betrachtet werden, die jetzt als bekannte Funktionen von \(x,\,x',\,x''\) usw. dargestellten Werte \(\xi',\xi'',\xi'''\) usw. annehmen. Dies lässt sich auch so ausdrücken: Soweit \(y,\,y',\,y''\) usw. als konstant und damit \(x',\,x'',\,x'''\) usw. nur als Funktionen der Variablen \(x\) betrachtet werden, müssen die folgenden \(n-1\) Differentialgleichungen erfüllt sein:
\[
\mathrm{d}x'=\xi'\mathrm{d}x,\qquad
\mathrm{d}x''=\xi''\mathrm{d}x,\qquad
\mathrm{d}x'''=\xi'''\mathrm{d}x,\quad\text{u.\ s.\ w.}
\]
oder, wenn man die ursprünglichen Gleichungen vorzieht, aus deren Kombination diese eigentlich entstanden sind, die folgenden:
\SourceLeaf{249}
\[
\begin{aligned}
&\quad *\qquad +\frac{(0,1)}{p}\mathbin{.}\mathrm{d}x'
 +\frac{(0,2)}{p}\mathbin{.}\mathrm{d}x''
 +\frac{(0,3)}{p}\mathbin{.}\mathrm{d}x'''+\text{usw.}\\
&=\frac{(1,0)}{p'}\mathbin{.}\mathrm{d}x\qquad *\quad
 +\frac{(1,2)}{p'}\mathbin{.}\mathrm{d}x''
 +\frac{(1,3)}{p'}\,\mathrm{d}x'''+\text{usw.}\\
&=\frac{(2,0)}{p''}\mathbin{.}\mathrm{d}x
 +\frac{(2,1)}{p''}\mathbin{.}\mathrm{d}x'\qquad *\quad
 +\frac{(2,3)}{p''}\mathbin{.}\mathrm{d}x'''+\text{usw.}\\
&=\text{u.\ s.\ w.}
\end{aligned}
\]
Die Integration dieser Gleichungen gehört jedoch in das Gebiet der gewöhnlichen Integralrechnung und wird hier vorausgesetzt. Sie enthält, allgemein gesprochen, \(n-1\) voneinander unabhängige Konstanten \(H,\,H',\,H'',\,H'''\) usw., die als gegebene Funktionen von \(x',\,x'',\,x'''\) usw. erscheinen, sodass man \(n-1\) endliche Gleichungen erhält:
\[
H=X,\qquad H'=X',\qquad H''=X'',\quad\text{u.\ s.\ w.}
\]
Dabei bezeichnen \(X,\,X',\,X''\) usw. diese Funktionen. Aus dieser Analyse folgt also, dass die vorgeschriebenen Bedingungen erfüllt sind, wenn man eben diese Funktionen für \(y,\,y',\,y''\) usw. wählt oder
\[
y=X,\qquad y'=X',\qquad y''=X'',\quad\text{u.\ s.\ w.}
\]
setzt.

\Needspace{32\baselineskip}
Eine Bemerkung wollen wir hier noch hinzufügen. Wir haben mit Bedacht angegeben, dass die Integration \emph{allgemein gesprochen} \(n-1\) voneinander unabhängige Konstanten liefert. In besonderen Fällen gilt dies nicht mehr, nämlich dann, wenn die Koeffizienten \(p,\,p',\,p''\) usw. so beschaffen sind, dass die obigen \(n-1\) Differentialgleichungen nicht voneinander unabhängig sind, sondern eine schon aus einer Kombination der übrigen abgeleitet werden kann. Hier ist auch die Bestimmung von \(\xi',\,\xi'',\,\xi'''\) usw. durch Elimination nicht mehr ausführbar. Dieser Fall müsste eigentlich als Ausnahme gesondert behandelt werden. Wir begnügen uns hier jedoch umso eher mit einer \emph{kurzen Andeutung}, wie man dabei vorgehen kann, da der Verfasser ihn nicht berührt hat. Um das Vorige anwenden zu können, braucht man nur diejenige der obigen Differentialgleichungen, die bereits in den anderen enthalten ist, durch irgendeine andere beliebige lineare Gleichung zwischen \(\mathrm{d}x,\,\mathrm{d}x',\,\mathrm{d}x''\) usw. zu ersetzen, die in diesen noch nicht enthalten ist. Am bequemsten ist es stets, eines dieser Differentiale \(=0\) zu setzen oder eine der Variablen \(x,\,x',\,x''\) usw.\SourceLeaf{250}
 als konstant zu behandeln, zum Beispiel \(x\), sofern nicht zufällig \(\mathrm{d}x=0\) schon aus den vorhandenen Gleichungen folgt. In diesem Fall ist eine der Integralgleichungen
\[
H=x
\]
und wenn man \(y=x\) setzt, lässt sich zeigen, dass sich im transformierten Ausdruck für \(\Omega\) stets von selbst \(q=0\) ergibt.

\begin{center}\rule{0.25\linewidth}{0.3pt}\end{center}


